% !TeX root = ../../tesis.tex
%%%%%%%%%%%%%%%%%%%%%%%%%%%%%%%%%%%%%%%%%%%%%%%%%%%%%%%%%%%%%%%%%%%
% SECCIÓN 5.5.2 — CENTRALIDAD DE INTERMEDIACIÓN B(v)
%%%%%%%%%%%%%%%%%%%%%%%%%%%%%%%%%%%%%%%%%%%%%%%%%%%%%%%%%%%%%%%%%%%
\subsection{Centralidad de Intermediación (\texorpdfstring{$B(v)$}{B(v)})}
\label{subsec:fase3-centralidad}

No todos los nodos de la red tienen el mismo peso en la conectividad global.
Algunos funcionan como puentes inevitables: la mayoría de las rutas mínimas
entre cualquier par de la red pasan por ellos, y su ausencia fragmentaría la
estructura o alargaría los tiempos de forma generalizada. La centralidad de
intermediación $B(v)$ mide exactamente esa condición: la proporción de rutas
mínimas del grafo que atraviesan cada nodo.

Dado un grafo $\mathcal{G}$ con nodos $\mathcal{V}$, la centralidad de
intermediación de un nodo $v$ se define como:

\begin{equation*}
    B(v) = \frac{1}{(N-1)(N-2)} \sum_{s \neq v \neq t} \frac{\sigma_{st}(v)}{\sigma_{st}}
    \tag{\ref{eq:betweenness}}
\end{equation*}

\noindent donde:
\begin{itemize}
    \item $\sigma_{st}$: cantidad de rutas mínimas entre los nodos $s$ y $t$.
    \item $\sigma_{st}(v)$: subconjunto de esas rutas que pasan por $v$.
    \item $N$: número de nodos del grafo. El factor $1/\bigl((N-1)(N-2)\bigr)$
    normaliza el valor para un grafo dirigido, de modo que $B(v) \in [0, 1]$.
\end{itemize}

El modelo calcula $B(v)$ con el algoritmo de Brandes, en la implementación de
NetworkX con la normalización activada. En grafos ponderados su complejidad es
$\mathcal{O}(VE + V^2 \log V)$; el cálculo se ejecuta sobre el componente
gigante (10{,}610~nodos).

Dicho de otro modo: hay estaciones que sostienen la red desde adentro sin que
los usuarios lo noten. No son necesariamente las más transitadas; son las que
aparecen en la ruta cuando no hay otra opción. Quitarlas no molesta a unos
cuantos pasajeros —desorganiza trayectos que atraviesan cuadrantes enteros de la
metrópoli.

El Cuadro~\ref{tab:bv-top10} presenta los diez nodos con mayor $B(v)$ calculados
sobre el componente gigante. El resultado más llamativo: los diez son estaciones
del Metro, y Tacubaya encabeza la lista con un valor de 0.2232.

\begin{table}[H]
    \centering
    \caption[Top 10 nodos por centralidad de intermediación $B(v)$]{%
    Diez nodos con mayor centralidad de intermediación ($B(v)$) en la
    red de transporte de la \zmvm{}. Corrida: 2026-10-04.}
    \label{tab:bv-top10}
    \small
    \begin{tabular}{|r|l|l|l|r|}
        \hline
        \textbf{\#} & \textbf{Estación} & \textbf{Sistema} &
        \textbf{Alcaldía} & \textbf{$B(v)$} \\
        \hline
         1 & Tacubaya                  & Metro & Miguel Hidalgo & 0.2232 \\
         2 & Mixcoac                   & Metro & Benito Juárez  & 0.1923 \\
         3 & Hidalgo                   & Metro & Cuauhtémoc     & 0.1515 \\
         4 & Balderas                  & Metro & Cuauhtémoc     & 0.1430 \\
         5 & Lázaro Cárdenas           & Metro & Cuauhtémoc     & 0.1365 \\
         6 & Zapata\textsuperscript{a} & Metro & Benito Juárez  & 0.1357 \\
         7 & Chilpancingo              & Metro & Cuauhtémoc     & 0.1348 \\
         8 & Zapata\textsuperscript{a} & Metro & Benito Juárez  & 0.1346 \\
         9 & Barranca del Muerto       & Metro & Álvaro Obregón & 0.1313 \\
        10 & Garibaldi y Lagunilla     & Metro & Cuauhtémoc     & 0.1303 \\
        \hline
    \end{tabular}
    \vspace{0.4em}

    \noindent\textsuperscript{a}\,Zapata aparece dos veces porque la
    estación se registra como dos nodos con coordenadas distintas.

    \fuentefigura{Elaboración propia con \texttt{VFTModel}
    \citep{cabrera2026VFTModel}. Corrida: 2026-10-04.}
\end{table}

El Metro monopoliza los diez primeros puestos pese a concentrar el 2.1\,\% de
las aristas de servicio real: su ventaja no es volumétrica sino posicional
---sus estaciones ocupan los puntos donde más rutas se cruzan.

La distribución geográfica del top~10 es igualmente estrecha: ocho de los diez
nodos quedan en Benito Juárez o Cuauhtémoc, lo que revela que la intermediación
de la red está concentrada en un corredor reducido del centro-poniente de la
cuenca.

El cuadro lee cada nodo por separado, pero una estación de correspondencia se
registra en el grafo como varios nodos, uno por línea o por sistema. La
Figura~\ref{fig:grafica-bv-nodos} ofrece la lectura complementaria: suma la $B(v)$ de
todos los nodos que comparten el nombre de la estación y apila la contribución
de cada sistema. El Cuadro~\ref{tab:bv-top10-agregado} resume los diez primeros
nombres. Con esta regla Tacubaya llega a 0.3279 y Chabacano, que no figura en el
cuadro por nodo, aparece en quinto lugar, porque sus nodos del Metro suman
0.2445. La agregación tiene dos límites que conviene declarar. Es por nombre
exacto y sin radio, de modo que también reúne paradas homónimas alejadas entre
sí: «Hidalgo» incluye paradas separadas hasta 17.6~km y «Chabacano», una parada
de la \textsc{rtp} a 15.5~km. Además, la fuente de la figura se limita a los
2{,}000 nodos con mayor $B(v)$; con los 10{,}610 nodos del componente gigante,
Hidalgo pasaría al primer lugar (0.3342). Por ello la figura se lee como un
indicador de la importancia de cada complejo de correspondencia y no sustituye
a la jerarquía por nodo.

\begin{table}[H]
    \centering
    \caption[Top 10 estaciones por suma de $B(v)$ por nombre]{%
    Diez nombres de estación con mayor suma de $B(v)$, con su desglose por
    sistema. Misma regla de agregación de la
    Figura~\ref{fig:grafica-bv-nodos}. Corrida: 2026-10-04.}
    \label{tab:bv-top10-agregado}
    \small
    \begin{tabular}{|r|l|r|l|}
        \hline
        \textbf{\#} & \textbf{Estación} & \textbf{Suma $B(v)$} &
        \textbf{Desglose por sistema} \\
        \hline
         1 & Tacubaya              & 0.3279 & Metro 0.3242; MB 0.0037 \\
         2 & Hidalgo               & 0.3265 & Metro 0.2710; Trolebús 0.0451; \textsc{rtp} 0.0105 \\
         3 & Mixcoac               & 0.2837 & Metro 0.2787; \textsc{rtp} 0.0051 \\
         4 & Zapata                & 0.2704 & Metro \\
         5 & Chabacano             & 0.2478 & Metro 0.2445; \textsc{rtp} 0.0033 \\
         6 & Barranca del Muerto   & 0.2431 & Metro 0.1313; \textsc{rtp} 0.0964; CC 0.0154 \\
         7 & Garibaldi y Lagunilla & 0.2364 & Metro \\
         8 & Balderas              & 0.2340 & Metro \\
         9 & Lázaro Cárdenas       & 0.1911 & Metro 0.1365; Trolebús 0.0546 \\
        10 & Bellas Artes          & 0.1812 & Metro 0.1503; MB 0.0171; Trolebús 0.0138 \\
        \hline
    \end{tabular}
    \fuentefigura{Elaboración propia con \texttt{VFTModel}
    \citep{cabrera2026VFTModel}.}
\end{table}

\begin{figure}[H]
    \centering
    \includegraphics[width=0.92\textwidth]{%
        Figures/Mapas/Tableau_Resources/RedActual/VFTModel/UniGrafica/%
        Grafica_VFTModel_2026_Centralidad_Intermediacion_Nodos}
    \caption[Centralidad de intermediación $B(v)$ --- suma por nombre de
    estación]{%
    Veinte nombres de estación con mayor suma de centralidad de intermediación
    ($B(v)$) en la red de transporte de la \zmvm{}. Cada barra suma la $B(v)$
    de los nodos que comparten el nombre exacto de la estación, apilada por
    sistema, a partir de los 2{,}000 nodos con mayor $B(v)$; por eso sus
    valores superan a los del Cuadro~\ref{tab:bv-top10}. Corrida: 2026-10-04.}
    \label{fig:grafica-bv-nodos}
    \fuentefigura{Elaboración propia con \textbf{Tableau};
    datos: \texttt{VFTModel} \citep{cabrera2026VFTModel}.}
\end{figure}

La centralidad topológica no siempre coincide con la afluencia real. Para
verificarlo, la Figura~\ref{fig:grafica-top10-estaciones} muestra las diez
estaciones con mayor volumen de usuarios según los registros de \texttt{Apímetro} ---un
contraste directo con la jerarquía de $B(v)$ del Cuadro~\ref{tab:bv-top10}.

\begin{figure}[H]
    \centering
    \includegraphics[width=0.92\textwidth]{%
        Figures/Mapas/Tableau_Resources/RedActual/Apimetro/UniGrafica/%
        Grafica_Apimetro_2026_Top_10_estaciones}
    \caption[Top 10 estaciones por afluencia de usuarios --- \zmvm{}]{%
    Diez estaciones con mayor volumen de usuarios registrado por
    \texttt{Apímetro} en la red de transporte de la \zmvm{}.
    Corrida: 2026-04-18.}
    \label{fig:grafica-top10-estaciones}
    \fuentefigura{Elaboración propia con \textbf{Tableau};
    datos: \texttt{Apímetro} \citep{cabrera2025apimetro}.}
\end{figure}

La Figura~\ref{fig:grafica-bv-sistema} agrega esos valores por sistema,
mostrando cuánto peso acumulado concentra cada modo en la intermediación global
de la red.

\begin{figure}[H]
    \centering
    \includegraphics[width=0.92\textwidth]{%
        Figures/Mapas/Tableau_Resources/RedActual/VFTModel/UniGrafica/%
        Grafica_VFTModel_2026_Centralidad_Intermediacion_sistema}
    \caption[Centralidad de intermediación $B(v)$ acumulada por sistema]{%
    Suma acumulada de $B(v)$ agrupada por sistema de transporte.
    El Metro concentra la mayor proporción de intermediación global
    de la red. Corrida: 2026-10-04.}
    \label{fig:grafica-bv-sistema}
    \fuentefigura{Elaboración propia con \textbf{Tableau};
    datos: \texttt{VFTModel} \citep{cabrera2026VFTModel}.}
\end{figure}

Para traducirlo al lenguaje urbano: imaginemos que cada par de colonias de la
\zmvm{} traza su ruta más rápida hacia el destino. Las estaciones que aparecen
con mayor frecuencia en ese ejercicio colectivo no son necesariamente las más
concurridas en términos de afluencia, sino las más difíciles de sustituir: su
ausencia obligaría a miles de trayectos a buscar desvíos más largos o, en el
peor caso, a no encontrarlos. $B(v)$ mide exactamente esa indispensabilidad
estructural.

La Figura~\ref{fig:mapa-betweenness} traslada esa jerarquía al territorio,
mostrando cómo la intermediación topológica se distribuye espacialmente sobre la
cuenca metropolitana.

\begin{figure}[H]
    \centering
    \includegraphics[width=0.92\textwidth]{%
        Figures/Mapas/Qgis_Resources/RedActual/Mapa_Betweenness}
    \caption[Mapa de centralidad de intermediación $B(v)$ --- \zmvm{}]{%
    Distribución espacial de la centralidad de intermediación ($B(v)$)
    en la red de transporte de la \zmvm{}. El corredor de alta
    intermediación traza una banda en el centro-poniente de la cuenca.
    Corrida: 2026-10-04.}
    \label{fig:mapa-betweenness}
    \fuentefigura{Elaboración propia con \textbf{QGIS};
    datos: \texttt{VFTModel} \citep{cabrera2026VFTModel}.}
\end{figure}

La concentración que revelan los datos tabulares adquiere forma geográfica en el
mapa: el corredor de alta intermediación no cubre la metrópoli de manera
uniforme, sino que traza una banda relativamente estrecha sobre el
centro-poniente, dejando amplias zonas periféricas con valores bajos o nulos.
Esas zonas ---las mismas que registraron menor cobertura y mayor Factor de
Desviación en Fase~1--- son también las que dependen de los nodos del corredor
central para cualquier viaje de largo aliento.

La Figura~\ref{fig:dashboard-bv} consolida la visión integral del indicador
sobre el conjunto de la red.

\begin{figure}[H]
    \centering
    \includegraphics[width=0.92\textwidth]{%
        Figures/Mapas/Tableau_Resources/RedActual/VFTModel/Dashboards/%
        Dashboard_VFTModel_2026_Centralidad_intermediacion}
    \caption[\termino{Dashboard} integral de centralidad de intermediación $B(v)$ --- \zmvm{}]{%
    Vista integral del indicador de centralidad de intermediación ($B(v)$)
    sobre la red de transporte de la \zmvm{}. La jerarquía usa la misma suma
    por nombre de estación de la Figura~\ref{fig:grafica-bv-nodos}.
    Corrida: 2026-10-04.}
    \label{fig:dashboard-bv}
    \fuentefigura{Elaboración propia con \textbf{Tableau};
    datos: \texttt{VFTModel} \citep{cabrera2026VFTModel}.}
\end{figure}

Los tres indicadores de Fase~3 ---tiempo de viaje, intermediación y
vulnerabilidad preliminar--- apuntan en la misma dirección: la red concentra su
capacidad de articulación en un corredor reducido del centro-poniente, mientras
la periferia sur y oriente queda estructuralmente subordinada a él. La sección
siguiente examina qué ocurre cuando alguno de esos nodos críticos desaparece de
la red.
